\documentclass[../../../include/open-logic-section]{subfiles}

\begin{document}

\olfileid[fr]{sth}{replacement}{extrinsic}
\olsection[Considérations extrinsèques]{Considérations extrinsèques sur le remplacement}

Commençons par examiner une tentative de justification
\emph{extrinsèque} du remplacement. Boolos en propose une :
\begin{quote}
  [\ldots] la raison d'adopter les axiomes de remplacement est fort
  simple : ils ont de nombreuses conséquences souhaitables et
  (apparemment) aucune conséquence indésirable. Outre les théorèmes
  concernant la conception itérative, ces conséquences comprennent
  une théorie satisfaisante, sinon idéale, des nombres infinis, ainsi
  qu'un résultat très souhaitable qui justifie les définitions
  inductives sur les relations bien fondées.
  \citep[229]{Boolos1971}
\end{quote}
L'idée essentielle de Boolos est que nous devrions justifier le
remplacement par ses fruits. Les fruits précis qu'il mentionne sont
ceux dont nous avons parlé dans les chapitres précédents. Le
remplacement nous a permis de montrer que les ordinaux de von Neumann
représentent remarquablement bien les types d'ordre des ensembles
bien ordonnés (c'est notre «~théorie satisfaisante, sinon idéale,
des nombres infinis~»).
Il nous a aussi permis de définir les $V_\alpha$, d'établir la notion
de rang et de démontrer le schéma de récurrence selon $\in$ (ce sont
nos «~théorèmes concernant la conception itérative~»). Enfin, le
remplacement nous permet de démontrer le théorème de récursion
transfinie (ce sont les «~définitions inductives sur les relations
bien fondées~»).

Il s'agit, en effet, de conséquences souhaitables. Mais suffisent-elles
à \emph{justifier} le remplacement ? \emph{Non}. Ou du moins, pas
directement.

Voici une première difficulté. Nous avons énuméré certaines conséquences
souhaitables du remplacement, mais nous aurions pu obtenir nombre
d'entre elles autrement. Ce fait est moins connu qu'il ne devrait
l'être ; prenons donc le temps d'expliquer la situation.

Il existe une théorie simple des ensembles, la théorie des niveaux
(\emph{Level Theory}), abrégée $\LT$.\footnote{Les premières versions
de $\LT$ sont proposées par \citet{Montague1965} et \citet{Scott1974} ;
\citet{Potter2004} l'a simplifiée et lui a consacré un livre ; plus
récemment, \citet{ButtonLT1} a encore simplifié $\LT$.}
Les axiomes de $\LT$ sont simplement l'extensionnalité, la séparation
et l'affirmation que tout ensemble est un sous-ensemble d'un certain
\emph{niveau}. La notion de «~niveau~» est habilement définie de
sorte que les niveaux se comportent comme nos bons vieux $V_\alpha$.
Les axiomes de $\LT$ se déduisent de $\ZF$. La théorie $\LT$ est
\emph{beaucoup} plus faible que $\ZF$. En effet, $\LT$ ne donne ni
l'axiome de la paire, ni celui
de l'ensemble des parties, ni celui de l'infini, ni celui du
remplacement. Appelons $\Zr$ la théorie obtenue en ajoutant l'infini
et l'ensemble des parties à $\LT$ ; on obtient alors aussi la paire,
si bien que $\Zr$ est au moins aussi forte que $\Z$. Mais $\Zr$ est en
fait strictement plus forte que $\Z$, car elle affirme en outre que
tout ensemble a un rang (d'où ma proposition de l'appeler $\Zr$).
De fait, $\Zr$ fournit une théorie des ordinaux parfaitement
satisfaisante, des résultats qui stratifient la hiérarchie en étapes
bien ordonnées, une démonstration de la récurrence selon $\in$ et
une \emph{version} de la récursion transfinie.

En somme, bien que Boolos ne l'ait pas su, toutes les conséquences
souhaitables qu'il mentionne pouvaient être obtenues \emph{sans}
le remplacement : il lui suffisait d'utiliser $\Zr$ plutôt que $\Z$.

(Étant donné cela, pourquoi avons-nous suivi la voie habituelle en
vous enseignant $\ZF$ plutôt que $\LT$ et $\Zr$ ? Pour deux raisons.
Premièrement, pour des raisons purement historiques, commencer par
$\LT$ est assez inhabituel ; nous voulions vous permettre de lire les
présentations plus classiques de la théorie des ensembles.
Deuxièmement, lorsque vous serez prêts à aborder $\LT$ et $\Zr$,
vous pourrez simplement lire \citealt{Potter2004} et
\citealt{ButtonLT1}.)

Bien entendu, puisque $\Zr$ est strictement plus faible que $\ZF$,
il existe des résultats que $\ZF$ démontre et que $\Zr$ laisse
indécidés. On pourrait donc tenter de justifier le remplacement
sur des bases extrinsèques en invoquant l'un de ces résultats.
Mais, dès que l'on sait utiliser $\Zr$, il devient assez difficile
de trouver beaucoup d'énoncés qui (a) soient établis par le
remplacement et non autrement, et (b) soient intuitivement vrais.
(Pour plus de détails, voir \citealt[\S13.2]{Potter2004}.)

L'essentiel est le suivant : pour fournir une justification
extrinsèque convaincante du remplacement, il faudrait trouver un
résultat que l'on \emph{ne peut pas} obtenir sans lui. Et ce n'est
pas une tâche facile.

Examinons une autre difficulté que rencontre toute tentative de
justification purement extrinsèque du remplacement. (Elle est
peut-être plus fondamentale que la première.) Boolos ne se contente
pas de dire que le remplacement a beaucoup de conséquences
souhaitables. Il affirme aussi qu'il n'a «~(apparemment) aucune
conséquence indésirable~». Mais la réserve «~apparemment~», entre
parenthèses, est assurément décisive.

Rappelons comment nous en sommes arrivés là : la compréhension
naïve a conduit à une contradiction, à laquelle nous avons répondu
en adoptant la conception cumulative et itérative des ensembles.
Cette conception s'accompagne d'un récit qui, nous l'espérons,
nous assure de sa cohérence. Mais si nous ne pouvons pas justifier
le remplacement à partir de ce récit, nous n'avons (pour l'instant)
aucune raison de croire que $\ZF$ est cohérente. Ou plutôt : nous
n'avons aucune raison de croire que $\ZF$ est cohérente, sinon le fait (peut-être
contingent) que personne n'a \emph{encore} découvert de contradiction.
C'est exactement en ce sens que la remarque de Boolos semble se
réduire à ceci : «~(apparemment) $\ZF$ est cohérente~». Nous devrions
exiger une garantie de cohérence plus solide.

Cette difficulté touche toute tentative \emph{purement} extrinsèque
de justifier le remplacement, c'est-à-dire toute justification
formulée uniquement en termes des conséquences (connues) de $\ZF$.
Nous voudrons donc chercher une justification \emph{intrinsèque}
du remplacement, qui fasse voir comment le récit que nous avons
fait des ensembles nous engage, d'une certaine manière, «~déjà~»
à accepter le remplacement.

\end{document}
